\section{研究设计与技术路线}

\subsection{总体研究框架}

\par 本研究的总体思路可概括为“特征驱动、类型分型、验证闭环”三个层次：
\begin{enumerate}
    \item 特征驱动：立足Solana架构特性，构建包含优先费用统计、新代币交互比例、指令级程序调用熵等Solana特有指标的行为特征体系，使检测模型能够捕捉Solana生态特有的行为模式。
    \item 类型分型：突破传统的二元检测框架，通过无监督聚类与监督学习相结合的策略，探索机器人的行为亚型（狙击型、套利型、做市型等），揭示不同类型机器人在Solana生态中的功能分化。
    \item 验证闭环：通过消融实验量化各特征组的贡献，通过典型案例分析将检测结果与链上实际行为相互印证，形成“特征→分类→验证”的完整闭环。
\end{enumerate}
\par 整体研究遵循“数据采集 → 特征工程 → 模型构建 → 验证分析”的四阶段技术路线。

\subsection{数据采集方案}

\par 候选地址的来源必须多元化，以减少单一来源的选择性偏差。我们将从以下四个渠道收集地址：
\begin{enumerate}
    \item 社区开源项目：GitHub上的Solana-MEV-Bot-Blackbook等列表由社区维护，收集了疑似MEV机器人的地址，具有一定可信度。可以直接整合该项目的最新列表。
    \item 区块链浏览器标签：Solscan和GMGN等浏览器允许用户为地址打标签，其中“机器人”“MEV”“狙击”等标签具有参考价值。可以通过API或手动方式采集这些地址。
    \item Dune Analytics：社区用户在Dune上创建了大量关于Solana交易的查询，部分查询专门追踪机器人活动。可以筛选相关查询的结果，提取地址。
    \item 反向追踪：从Meme币发行事件入手，追踪新代币上线后前几笔交易的参与者。如果某个地址频繁出现在不同新代币的首批交易中，其机器人的嫌疑就很大。
\end{enumerate}
\par 采集到的候选地址将经过初步清洗，并执行活跃度筛选。筛选标准设定为：在2026年2月内存在至少一个连续7天的活跃窗口，且该窗口内交易数不低于1800笔（日均约257笔）。这一高门槛旨在过滤掉普通用户的人工交易，锁定高频自动化程序。
\par 交易数据采集通过Solana RPC接口实现。对于每个通过筛选的地址，将获取其在2月内的所有交易签名，然后逐个获取交易详情。采集过程将采用多节点轮询策略和线程池控制，避免触发RPC限流。

\subsection{凝练Solana特有特征及其行为含义}

\par 特征设计的核心思想是：从多个维度刻画地址的行为模式，使机器人与普通用户的差异能够被量化。本研究将特征分为五个类别，共计20余维。


\begin{table}[htbp]
\centering
\caption{Solana机器人检测特征体系}
\label{tab:feature_system}
\begin{tabular}{p{2.3cm} p{3.2cm} p{7cm}}
\hline
\textbf{特征类别} & \textbf{特征名称} & \textbf{特征含义与设计动机} \\
\hline
\multirow{4}{*}{\textbf{基础活跃度}} & 总交易数 & \multirow{4}{7cm}{反映地址的整体交易强度，高频是机器人的基本特征。} \\
& 日均交易数 & \\
& 单日最大交易数 & \\
& 活跃天数 & \\
\hline
\multirow{5}{*}{\textbf{时间规律}} & 夜间交易比例 & \multirow{5}{7cm}{机器人通常全天活跃、间隔均匀，而人类用户受作息影响。} \\
& 周末交易比例 & \\
& 平均交易间隔 & \\
& 间隔变异系数 & \\
& 小时交易熵 & \\
\hline
\multirow{6}{*}{\textbf{交互多样性}} & 调用程序ID数量 & \multirow{6}{7cm}{程序调用熵衡量交互的集中度；新代币交互比例捕捉机器人对Meme币生命周期的参与程度。} \\
& 程序调用熵 & \\
& 交互代币种类数 & \\
& 代币交互熵 & \\
& 与核心DEX交互占比 & \\
& 新代币交互比例 & \\
\hline
\multirow{3}{*}{\textbf{Gas策略}} & 平均优先费用 & \multirow{3}{7cm}{机器人为抢跑愿支付高额优先费用，且费用波动反映其竞价策略的激进程度。Solana的本地费用市场使这一特征更具区分度。} \\
& 优先费用变异系数 & \\
& 优先费用与网络拥堵相关性 & \\
\hline
\multirow{3}{*}{\textbf{指令级特征}} & 平均指令数 & \multirow{3}{7cm}{Solana交易包含多指令，机器人常通过复杂指令序列实现套利策略，解析指令级特征可揭示行为复杂性。} \\
& 指令类型多样性 & \\
& 程序调用模式 & \\
\hline
\end{tabular}
\end{table}
\par 所有特征将在计算前进行归一化处理，以消除量纲影响。特征计算代码将封装为独立模块，支持对新地址一键生成特征向量。

\subsection{检测模型设计}
\par 在完成特征工程的基础上，本研究的核心是构建一套能够有效识别并分型Solana MEV机器人的检测模型。为实现这一目标，我们设计了从问题定义到模型验证的完整研究路径，重点围绕“Solana特有行为模式捕捉”、“机器人行为亚型分治”两个维度展开。
\par 首先，明确研究的核心问题——如何利用链上交易特征检测Solana MEV机器人。在此基础上，立足Solana架构，构建包含优先费用竞价痕迹、新代币生命周期参与度、指令级程序调用结构等特有指标的特征体系，将抽象的机器人行为转化为可量化的特征向量。
\par 其次打破传统的“机器人/非机器人”二元分类框架。先通过无监督学习探索数据的自然分布，识别出狙击、套利、做市等可能存在的机器人行为亚型。在此基础上，构建监督学习模型实现对地址的精确分类，并探索多分类模型以区分不同亚型。
\par 通过消融实验，量化Solana特有特征组对模型性能的贡献，以证明其必要性。
\par 在常规模型验证的基础上，探索引入LLM 进行辅助标注与行为模式解读，利用其语义理解能力提升数据标注效率和对异常地址的聚类解释能力，为检测方法提供新的技术路径。
\par 基于上述研究思路，我们设计并实施以下模型构建方案：
\begin{itemize}
    \item 基线模型：构建一套基于启发式规则的检测方法作为性能基线。规则设计参考现有研究，并结合Solana特性，如设定高夜间交易比例、高优先费用、低程序调用多样性等阈值。此步骤旨在为后续机器学习模型提供一个可解释的性能下限。
    \item 监督学习：采用集成学习算法（随机森林、XGBoost、LightGBM）作为核心分类器。我们将利用人工标注的机器人及普通用户地址构建训练集。模型训练后，不仅评估其准确率、F1分数等指标，更将特征重要性分析作为关键一环，以验证优先费用、新代币交互比例等Solana特有特征是否在分类中起到了决定性作用。
    \item 无监督学习与行为分型：为弥补标签数据的稀缺性并探索机器人亚型，我们将应用高斯混合模型和DBSCAN等聚类算法对未标注地址进行行为分型。通过分析聚类结果中各簇的特征均值，我们能够识别出不同的行为模式。
    \item 细粒度多分类：若聚类结果能清晰定义出不同的机器人亚型，我们将基于此构建多分类标注集，训练能够区分狙击、套利等不同机器人的多分类模型，实现更精细的检测与识别。
\end{itemize}



\section{前期实验基础}

\subsection{复现论文}

\par 在正式进入Solana机器人检测研究之前，我对《Detecting Financial Bots on the Ethereum Blockchain》进行了完整的代码复现。通过亲手走通流程，我大致了解了其方法论精髓，并发现其中可能隐藏的坑。
\par 复现过程让我对几个关键问题有了切身体会。当我使用论文提供的完整数据集时，随机森林的准确率能达到83\%，聚类纯度可达82.6\%；但一旦将数据缩减到十分之一，模型性能就出现明显下滑，召回率从0.77跌至0.50。这说明在Solana上构建数据集时，必须保证足够的地址数量，尤其是让各类机器人的样本都达到一定规模，否则模型很容易过拟合或欠拟合。
\par 类别平衡同样是绕不开的问题。原论文的多类MEV分类任务中，清算（Liquidation）类样本有111个，模型能取得93\%的召回率；而我们在复现时该类别仅剩3个样本，召回率直接跌到33\%。这说明后续若要做细粒度分类，必须确保每个子类有足够的代表。
\par 特征完整性的影响也不容忽视。原论文基于完整时间窗口（15天）计算了83维特征，复现时由于数据截断，部分依赖长期行为模式的特征（如周期性休眠）无法准确计算，削弱了特征的表达能力。这提示在Solana数据采集中，必须保证时间窗口的完整性，不能为了节省RPC请求而截断历史。
\par 同时我也看清了启发式规则的局限：组合规则的召回率不足44\%，远低于机器学习模型。这坚定了以机器学习为核心方法的技术路线，同时也认可启发式规则作为可解释基线的价值。

\subsection{前期实验的积累}
\par 在复现论文的启发下，我针对Solana开展了前期探索性实验。实验目标很明确：验证从Solana链上采集数据、计算行为特征、筛选高频地址的整个流程是否可行。
\par 数据来源方面，我从GMGN上手动收集了一些交易量巨大、胜率或盈利率异常高的地址，同时整合了GitHub上的Solana-MEV-Bot-Blackbook项目。最终得到80个候选地址。
\par 采集策略上，我设定了两种模式：default模式从当前时间向前回溯寻找符合条件的窗口，feb模式严格限定在2026年2月内。通过滑动窗口检测（窗口长度7天，交易量门槛1800笔），最终有56个地址通过筛选。采集过程中验证了动态回溯机制的有效性：程序会持续获取签名直到覆盖整个2月，确保不会因FETCH\_LIMIT截断而遗漏早期交易。
\par 在特征计算方面实现了12维特征，涵盖基础活跃度（日均交易数、单日最大交易数）、时间规律（夜间比例、周末比例、平均间隔、小时熵）和交互多样性（程序种类、代币种类、程序熵、代币熵）。所有特征基于2月内的交易时间戳和交易详情计算得出，单地址最多取500笔交易详情用于程序/代币解析，以平衡精度和RPC负载。
\par 初步结果让人眼前一亮。日均交易数分布呈左偏，多数地址集中在1000~3000笔/天，这种高强度交易绝非人工可为。小时熵分布在3.6至4.65之间，覆盖了从相对集中到接近均匀的宽范围，暗示机器人行为存在分化——有的全天均匀活跃，有的集中在特定时段。平均交易间隔多数在50秒以内，进一步印证了高频特征。程序调用种类集中在5种左右（Jupiter、Raydium为主），体现了Solana机器人的专业化；代币交互种类则呈现两极分化，多数地址交易代币较少，少数地址涉及大量代币，可能对应做市或多策略角色。